% solids_stressstrain: stress-strain curves (Vol II Ch 1, fig:solids-stressstrain)
% Schematic, strain axis NOT to scale: for a real metal the elastic region ends near a strain
% of 0.002 while fracture comes at 0.2-0.5, so a to-scale plot would hide the straight part.
% Left, a ductile metal, in units where the proportional limit is at (0.08, 0.8):
%   O-A   linear, sigma = 10 eps                                     (Hooke's law)
%   A-B   sigma = 0.8 + 10(eps-0.08) - 62.5(eps-0.08)^2, eps in [0.08,0.12] (elastic, not linear)
%   B-D   sigma = 1.1 + 0.45(1 - exp(-(eps-0.12)/0.09)), eps in [0.12,0.75] (plastic, hardening)
%         (slope 5 at B, matching A-B there)
%   D-E   sigma = sigma_D - 3(eps-0.75)^2, eps in [0.75,0.95]          (necking to fracture)
%   Unloading from C (eps = 0.40) follows a line of the elastic slope 10 back to
%   eps = 0.40 - sigma_C/10: the strain left over is the permanent set.
% Right, rubber (an elastomer), schematic: loading sigma = 0.3(lam - 1/lam^2) + 0.004 eps^4
%   with lam = 1 + eps (neo-Hookean, stiffening at large stretch), unloading lower,
%   sigma_load * (eps/6)^0.5.  No straight region; the loop's area is energy lost as heat.
\begin{tikzpicture}
  % ------------------------------------------------------------ metal
  \begin{axis}[bookplot,
      width=7.0cm, height=4.6cm, scale only axis,
      xmin=0, xmax=1.08, ymin=0, ymax=1.85,
      xlabel={strain $\epsilon$},
      xlabel style={at={(axis description cs:1,0)},anchor=north east,yshift=-2pt},
      ylabel={stress $\sigma$},
      ylabel style={at={(axis description cs:0,1)},anchor=south,rotate=-90}]
    \pgfmathsetmacro\sD{1.1+0.45*(1-exp(-(0.75-0.12)/0.09))}
    \pgfmathsetmacro\sC{1.1+0.45*(1-exp(-(0.40-0.12)/0.09))}
    \pgfmathsetmacro\eS{0.40-\sC/10}
    % unloading from C and the permanent set
    \addplot[guide,domain=0:1,samples=2] ({0.40-\sC/10*(1-x)},{\sC*x});
    \draw[dim] (axis cs:0,0.12) -- (axis cs:\eS,0.12);
    \node[anchor=north,font=\footnotesize,text=cSteel] at (axis cs:{\eS/2},0) {permanent set};
    % the curve
    \addplot[cPrim,domain=0:0.08,samples=2] {10*x};
    \addplot[cPrim,domain=0.08:0.12,samples=20] {0.8+10*(x-0.08)-62.5*(x-0.08)^2};
    \addplot[cPrim,domain=0.12:0.75,samples=80] {1.1+0.45*(1-exp(-(x-0.12)/0.09))};
    \addplot[cPrim,domain=0.75:0.95,samples=30] {\sD-3*(x-0.75)^2};
    % landmarks
    \node[dot] at (axis cs:0.08,0.8) {};
    \node[anchor=west] at (axis cs:0.09,0.74) {A};
    \node[dot] at (axis cs:0.12,1.1) {};
    \node[anchor=south east] at (axis cs:0.125,1.12) {B};
    \node[dot] at (axis cs:0.40,\sC) {};
    \node[anchor=south] at (axis cs:0.40,\sC) {C};
    \node[dot] at (axis cs:0.75,\sD) {};
    \node[anchor=south] at (axis cs:0.75,\sD) {D};
    \node[cCrim,font=\normalsize] at (axis cs:0.95,{\sD-0.12}) {$\times$};
    \node[anchor=west] at (axis cs:0.965,{\sD-0.12}) {E};
    \node[anchor=north west,align=left,font=\footnotesize,text=cSteel] at (axis cs:0.47,0.98)
      {A proportional limit\\ B elastic limit, yield\\ C unloaded from here\\ D ultimate stress\\ E fracture};
  \end{axis}
  % ------------------------------------------------------------ rubber
  \begin{scope}[xshift=9.0cm]
  \begin{axis}[bookplot,
      width=5.4cm, height=4.6cm, scale only axis,
      xmin=0, xmax=6.6, ymin=0, ymax=8.4,
      xlabel={strain $\epsilon$},
      xlabel style={at={(axis description cs:1,0)},anchor=north east,yshift=-2pt},
      ylabel={stress $\sigma$},
      ylabel style={at={(axis description cs:0,1)},anchor=south,rotate=-90}]
    \addplot[cPrim,domain=0:6,samples=80] {0.3*((1+x)-1/(1+x)^2)+0.004*x^4};
    \addplot[cTeal,domain=0:6,samples=80] {(0.3*((1+x)-1/(1+x)^2)+0.004*x^4)*sqrt(x/6)};
    \node[cPrim,anchor=east] at (axis cs:4.3,4.0) {loading};
    \node[cTeal,anchor=west] at (axis cs:4.6,1.9) {unloading};
    \node[anchor=north west,align=left,font=\footnotesize,text=cSteel] at (axis cs:0.3,7.9)
      {rubber: no straight part;\\ strains of several hundred\\ per cent; the loop's area\\ is energy lost as heat};
  \end{axis}
  \end{scope}
\end{tikzpicture}
